\Needspace{18\baselineskip}
\section{Associativity, units, and exact context composition}
\label{sec:composition}

We now compose addressed substitutions, rather than merely compare renamed
versions of one substitution. If $s$ selects an occurrence in $G$, inserting
$H$ there creates a copy of every occurrence of $H$. A site $t$ in $H$ thus
induces a site in $\mathsf{plugAt}(s,H)$. The main theorem identifies substitution
at that induced site with first substituting within $H$ and then inserting the
result at $s$. The identification is a concrete recursive isomorphism; on
isomorphism classes it is equality. We also prove the corolla unit laws and
an exact context-validity criterion for the composite, including cases where
an intermediate collision is removed by the final substitution.

\subsection{Shallow associativity with explicit carrier bijections}

Let $C$ be a graph, $e$ a selected edge, $H$ a replacement of its interface,
$f$ an edge of $H$, and $A$ a replacement of the interface of $f$. Compare
\[
 L=(C[e:=H])[\mathsf{inr}(f):=A],\qquad
 R=C[e:=H[f:=A]].
\]
These have the same outer interface, but different carrier records.

\begin{definition}[Reassociation bijections]\label{term:reassociation}
The internal-vertex bijection from $L$ to $R$ is
\[
 (C_V+H_V)+A_V\longrightarrow C_V+(H_V+A_V),
\]
which keeps each original name and changes only its sum tags. Its three cases
send a host vertex to the host summand, a middle vertex to the middle summand,
and a final vertex to the final summand. The edge bijection does the same for
surviving host edges, surviving middle edges, and final inserted edges.
\end{definition}
More explicitly, the left surviving-edge carrier consists of edges of
$(C_E\setminus\{e\})+H_E$ other than $\mathsf{inr}(f)$. A surviving host
edge maps to the host summand on the right. An edge $\mathsf{inr}(y)$ maps to
the surviving-middle summand, with $y\ne f$ obtained by injectivity of the sum
constructor. Final edges map to final edges. The inverse maps reverse these
three cases; inverse equations are proved by cases, with proof irrelevance
identifying the subtype certificates.
\lean{Graph/PlugAssociativity}{13}{18}{Net.PlugAssoc.vertices}
\lean{Graph/PlugAssociativity}{21}{32}{Net.PlugAssoc.edges}

\begin{lemma}[Attachment and port commutation]
The source and target positions of the inserted edge $\mathsf{inr}(f)$ are the
images, under insertion into $C$, of the corresponding positions of $f$ in $H$.
The reassociation port map commutes with inclusion of original host ports,
inclusion of middle ports, and insertion of final replacement ports.
\end{lemma}
\begin{proof}
The attachment assertions unfold positional lookup in the inserted incidence
lists and use lookup through a list map. For host ports, case analysis on
internal, input, and output ports proves the equation directly. For middle
ports, internal vertices are immediate; an input or output is attached to a
host port, so use the host equation. For final ports, internal vertices are
again immediate. A boundary port attaches at $f$; rewrite that attachment by
the positional lookup assertion and apply the middle-port equation.
No argument assumes that an input attachment differs from an output attachment.
\end{proof}
\lean{Graph/PlugAssociativity}{35}{37}{Net.PlugAssoc.sourcePort_inserted}
\lean{Graph/PlugAssociativity}{40}{42}{Net.PlugAssoc.targetPort_inserted}
\lean{Graph/PlugAssociativity}{45}{48}{Net.PlugAssoc.map_host}
\lean{Graph/PlugAssociativity}{51}{58}{Net.PlugAssoc.map_middle}
\lean{Graph/PlugAssociativity}{61}{75}{Net.PlugAssoc.map_insert}

\begin{theorem}[Shallow associativity]
There is a literal-label graph isomorphism $L\cong R$.
\end{theorem}
\begin{proof}
Use the two reassociation bijections. Vertex and edge labels agree by the
three carrier cases. For sources, split an edge into surviving host, surviving
middle, or final replacement. In the host case the left source list maps the
original host list through two inclusions and reassociation; combine list maps
and apply host commutation. The middle case uses middle commutation, and the
final case uses final insertion commutation. These give exactly the three
right-hand source lists. The target proof is identical with target lists.
The isomorphism preserves literal labels and is valid for self-loops and cycles.
\end{proof}
\lean{Graph/PlugAssociativity}{80}{118}{Net.plug_assoc}

\subsection{Composition at arbitrary box depth}

\begin{lemma}[Overwrite and insertion of box bodies]
Replacing the same box body twice has exactly the effect of the last replacement.
Changing the body of an inserted box commutes, up to literal-label isomorphism,
with changing that body before shallow insertion.
\end{lemma}
\begin{proof}
For overwrite, the two graphs have the same carriers, incidence, and internal
labels. At the selected edge both final labels contain the last body. At any
other edge both conditional updates keep the original label. Extensionality of
the label function and proof irrelevance give equality of the graph records.

For commutation with insertion use identity carrier bijections. Surviving host
labels are unchanged. On inserted edges the two tests for being the selected
edge agree, because $\mathsf{inr}(x)=\mathsf{inr}(f)$ exactly when $x=f$.
Incidence is unchanged by box replacement. The insertion functions also agree:
internal ports have the same tags, and boundary ports attach to the same outer
host positions. List mapping then proves both incidence equations.
\end{proof}
\lean{Graph/DeepAssociativity}{10}{20}{Linear.replaceBox_overwrite}
\lean{Graph/DeepAssociativity}{23}{62}{Linear.plug_replaceBox_iso}

\begin{definition}[Composed site]\label{term:composed-site}
Let $s:\mathsf{Site}(G,r,J,K)$ and $t:\mathsf{Site}(H,q,U,V)$, where $H$ has
the replacement interface of $s$. Define
\[
 s\star t:\mathsf{Site}(\mathsf{plugAt}(s,H),q,U,V).
\]
If $s$ is shallow and $t$ is shallow, select the inserted edge
$\mathsf{inr}(f)$. If $s$ is shallow and $t$ enters a box, enter the corresponding
inserted box and retain $t$'s inner address. If $s$ enters a box, rebuild that
box with the first replacement already performed, then descend to the recursively
composed inner site.
\end{definition}
This is recursion on the outer address. The new enclosing box carries exactly
the body produced by the first substitution, so the box-label equation is the
selected branch of the update. Local depth and interface indices are preserved
by construction. Although the surrounding graph operations use classical
finite-subset constructions, no isomorphism or target address is selected by
an existence theorem here.
\lean{Graph/DeepAssociativity}{66}{77}{Linear.Site.comp}

\begin{theorem}[Deep associativity]
For every compatible $s,t,A$,
\[
 \mathsf{plugAt}(s\star t,A)\cong_d
 \mathsf{plugAt}(s,\mathsf{plugAt}(t,A)).
\]
No uniqueness, freshness, or non-self-loop premise is required.
\end{theorem}
\begin{proof}
Induct on $s$. When $s$ and $t$ are shallow use the explicit shallow associator
and embed its literal-label isomorphism into recursive isomorphism. When $s$
is shallow and $t$ enters a box, use the box/insertion commutation lemma.
When $s$ enters a box, the left expression first updates that box to contain
$\mathsf{plugAt}(s_{\rm inner},H)$ and later updates it again after following
the composed inner address. Overwrite reduces these two updates to one. The
induction hypothesis compares its final body to the body obtained by performing
the inner substitution before the outer one. Box replacement congruence lifts
that body isomorphism through the enclosing box.
\end{proof}
\lean{Graph/DeepAssociativity}{80}{94}{Linear.plugAt_assoc}

\begin{lemma}[Quotient composition and three-stage coherence]
On recursive isomorphism classes,
\[
 P_{s\star t}([A])=P_s(P_t([A])),\qquad
 P_s=\mathsf{plugAtModuloIso}(s,-).
\]
For a third compatible site $u$, the two functions associated with
$(s\star t)\star u$ and $s\star(t\star u)$ also agree on every replacement class.
\end{lemma}
\begin{proof}
For the binary equation use quotient induction and quotient soundness applied
to deep associativity on representatives. For three stages, apply the binary
equation twice to each side; both reduce to $P_s(P_t(P_u([A])))$.
The ambient concrete hosts of the two composite addresses can differ, so this
is a statement about the induced quotient results, not literal equality of
address records in different hosts.
\end{proof}
\lean{Graph/DeepAssociativity}{97}{102}{Linear.plugAtModuloIso_comp}
\lean{Graph/DeepAssociativity}{205}{210}{Linear.plugAtModuloIso_comp_three}

The selected-hole observation of $s\star t$ equals that of $t$. Induct on $s$;
for shallow $s$, either form of $t$ retains the same selected label, and for
inside $s$ the definition discards the enclosing box and the induction hypothesis
applies.
\lean{Graph/DeepAssociativity}{105}{110}{Linear.Site.selectedHole_comp}

\subsection{Exact inventories and context validity of composites}

Write $R_s(h)$ for $\mathsf{SiteRetainedHole}(s,h)$ and $H_A(h)$ for
$\mathsf{HasHole}(A,h)$.

\begin{lemma}[Retained inventories compose by union]
Removing an inserted shallow edge $\mathsf{inr}(f)$ retains exactly the holes
under surviving original-host edges and surviving middle-graph edges. More
generally,
\[
 R_{s\star t}(h)\quad\Longleftrightarrow\quad R_s(h)\lor R_t(h).
\]
Consequently the final hole inventory is exactly
$R_s(h)\lor R_t(h)\lor H_A(h)$.
\end{lemma}
\begin{proof}
For shallow inserted-edge removal, decompose a retained occurrence by its outer
edge tag. A surviving host tag already supplies its inequality with the original
selected edge. An inserted tag survives precisely when its middle edge differs
from $f$. Its fibre position and hole observation are unchanged. Conversely,
construct the appropriate tagged edge from either of these witnesses.

For the general statement induct on $s$. A shallow $s$ and shallow $t$ give
the preceding result directly. If $t$ enters a box, append its retained inner
inventory and reassociate disjunction. For inside $s$, replacement of its
selected box preserves all retained outer-edge fibres; apply that fact and
the induction hypothesis to the inner composite, then reassociate disjunction.
Finally apply the exact inventory theorem for replacement at $s\star t$.
\end{proof}
\lean{Graph/DeepAssociativity}{114}{124}{Linear.retainedHole_inserted}
\lean{Graph/DeepAssociativity}{127}{143}{Linear.siteRetainedHole_comp}
\lean{Graph/DeepAssociativity}{147}{153}{Linear.plugAt_comp_hasHole_iff}

\begin{lemma}[Composite freshness]
Exact freshness for $(s\star t,A)$ is equivalent to both
\[
 \forall h,\ \neg(R_s(h)\land H_A(h))
 \quad\text{and}\quad
 \mathsf{SiteFresh}(t,A).
\]
\end{lemma}
\begin{proof}
Use the retained-union theorem. In the forward direction include either retained
inventory into the union and apply composite freshness. In the reverse direction,
split a retained witness according to the union and apply the corresponding
freshness assumption.
\end{proof}
\lean{Graph/DeepAssociativity}{156}{168}{Linear.siteFresh_comp_iff}

\begin{theorem}[Sharp context closure for composition]
Suppose $G$ and $H$ have unique holes. Then
$\mathsf{plugAt}(s\star t,A)$ has unique holes if and only if all three conditions
hold:
\begin{enumerate}
\item $A$ has unique holes;
\item $A$ is fresh for the composite site $s\star t$;
\item $\forall h,\ \neg(R_s(h)\land R_t(h))$.
\end{enumerate}
Thus the three inventories $R_s,R_t,H_A$ must be pairwise disjoint.
The intermediate graph $\mathsf{plugAt}(s,H)$ need not have unique holes.
\end{theorem}
\begin{proof}
By deep associativity it suffices to check uniqueness of
$\mathsf{plugAt}(s,\mathsf{plugAt}(t,A))$. Apply sharp context closure first
with host $G$, then with host $H$. This gives uniqueness of $A$, freshness for
$(t,A)$, and disjointness of $R_s$ from the hole inventory of
$\mathsf{plugAt}(t,A)$. The latter inventory is $R_t\cup H_A$ by the exact
inventory theorem. Splitting that union gives disjointness of $R_s$ from $R_t$
and from $H_A$. Combine the last condition with freshness for $(t,A)$ using
the composite-freshness lemma. Conversely, reverse each implication: the
pairwise exclusions make the inner result fresh for $s$, and both applications
of sharp closure establish the final result's uniqueness.
\end{proof}
\lean{Graph/DeepAssociativity}{173}{191}{Linear.plugAt_comp_unique_iff}

\begin{definition}[Certified composed plugging]\label{term:certified-composition}
Given contexts $G,H,A$, compatible sites $s,t$, composite freshness, and
disjoint retained inventories $R_s,R_t$, package
$\mathsf{plugAt}(s\star t,A)$ with the uniqueness proof just obtained.
This defines \texttt{Context.plugAtComp}; it requires no certificate that the
intermediate plugged graph is a context.
\end{definition}
\lean{Graph/DeepAssociativity}{195}{201}{Linear.Context.plugAtComp}

\subsection{Corolla units}

\begin{definition}[Substitution unit]\label{term:substitution-unit}
For a signature label $o$, let $\mathbf c_o$ be its one-edge corolla, with its
sole edge selected. This is a substitution unit of interface
$(\sigma_{\rm in}(o),\sigma_{\rm out}(o))$. It is an addressed graph, with no uniqueness premise on its label;
no free wire identifying an input vertex with an output vertex is introduced.
\end{definition}
We prove $\mathbf c_o[0:=H]\cong H$ and
$C[e:=\mathbf c_{\ell(e)}]\cong C$ using literal-label isomorphisms.

A corolla has no surviving edge after removing zero: an element of
$\{i:\mathsf{Fin}(1)\mid i\ne0\}$ contradicts the fact that every index below
one is zero. Its $i$th source and target attachment positions are its input
and output ports at $i$, respectively. These equations follow by unfolding
lookup in the mapped finite-range lists.
\lean{Graph/PlugUnits}{10}{10}{Net.PlugUnit.noEdge}
\lean{Graph/PlugUnits}{13}{14}{Net.PlugUnit.corolla_sourcePort}
\lean{Graph/PlugUnits}{17}{18}{Net.PlugUnit.corolla_targetPort}

\begin{lemma}[Left unit]
Replacing the sole edge of $\mathbf c_o$ by $H$ yields a graph isomorphic to $H$.
\end{lemma}
\begin{proof}
The vertex bijection discards the empty left summand in $\varnothing+H_V$.
The edge bijection discards the impossible surviving-corolla summand and retains
$H_E$. Their inverses include into the right summands. Both inverse laws follow
by cases, eliminating impossible elements.

Mapping an inserted port through the vertex projection gives the original port:
internal ports are immediate; input and output ports use the corolla attachment
equations. For each inserted edge, combine the incidence list maps and apply
this port equation, leaving the original source or target list. Surviving
corolla edges are impossible. All labels agree directly on the remaining
summands.
\end{proof}
\lean{Graph/PlugUnits}{21}{25}{Net.PlugUnit.leftVertices}
\lean{Graph/PlugUnits}{28}{32}{Net.PlugUnit.leftEdges}
\lean{Graph/PlugUnits}{35}{40}{Net.PlugUnit.left_insert}
\lean{Graph/PlugUnits}{88}{111}{Net.corolla_plug_iso}

\begin{lemma}[Right unit]
Replacing $e$ in $C$ by its own labelled corolla yields a graph isomorphic to $C$.
\end{lemma}
\begin{proof}
The vertex bijection discards the empty right summand. The edge bijection maps
each surviving host edge to itself and the inserted corolla edge to $e$.
Its inverse tests whether an edge equals $e$: if so use the sole corolla edge;
otherwise use the surviving-host subtype. The inverse laws follow from this
test and uniqueness of an element of $\mathsf{Fin}(1)$.

The vertex projection cancels inclusion of host ports. Enumerating all source
positions of $e$ reconstructs $C$'s source list: the lengths agree by typing,
and each indexed element agrees by lookup in the finite-range map. The target
list is reconstructed in the same way. Thus surviving-edge incidence agrees
by cancellation of host inclusion; the inserted corolla's full boundary list
maps to exactly the source or target list of $e$. Labels agree on both edge
cases. This proof also covers an edge whose source and target share vertices.
\end{proof}
\lean{Graph/PlugUnits}{43}{47}{Net.PlugUnit.rightVertices}
\lean{Graph/PlugUnits}{50}{63}{Net.PlugUnit.rightEdges}
\lean{Graph/PlugUnits}{66}{68}{Net.PlugUnit.right_host}
\lean{Graph/PlugUnits}{71}{76}{Net.PlugUnit.sourcePort_list}
\lean{Graph/PlugUnits}{79}{84}{Net.PlugUnit.targetPort_list}
\lean{Graph/PlugUnits}{114}{145}{Net.plug_corolla_iso}

\Needspace{16\baselineskip}
\subsection{Restoration and unit equations at deep sites}

\begin{definition}[Original selected fragment]\label{term:original-fragment}
The original fragment at a site $s$ is the corolla carrying the label of its
selected edge. For an inside site descend to its inner site. If the selected
edge is itself a box, its original fragment includes that complete box label
and its original body. The fragment has exactly the local depth and interface
required by $s$.
\end{definition}
\lean{Graph/DeepUnits}{11}{15}{Linear.Site.original}

\begin{lemma}[Deep restoration]
Substituting the original fragment at any site restores the entire host up to
recursive isomorphism.
\end{lemma}
\begin{proof}
First note that reinstalling a box's original body gives a literal-label
isomorphism with identity carriers: at the selected edge use its original
box-label equality, and elsewhere no label changes. Incidence and internal
labels are untouched.

Now induct on the site. For a shallow site apply the right unit. For an inside
site, the induction hypothesis restores its inner host up to recursive
isomorphism. Box congruence transports that isomorphism to the surrounding
updated box. Compose it with the just-proved isomorphism for reinstalling the
original body.
\end{proof}
\lean{Graph/DeepUnits}{18}{33}{Linear.replaceBox_original_iso}
\lean{Graph/DeepUnits}{37}{43}{Linear.plugAt_original}

Consequently, plugging the original fragment class at $s$ returns the original
host class, by quotient soundness. A selected one-edge corolla acts as identity
on every replacement class of its interface: quotient induction reduces this
to the left-unit isomorphism. Together with the composition equations above,
these are unit and associativity laws for the induced substitution operations
on recursive isomorphism classes. The quotient of pointed hosts and the
independent-site interchange law are developed in Section~\ref{sec:deep-interchange}.
\lean{Graph/DeepUnits}{46}{49}{Linear.plugAtModuloIso_original}
\lean{Graph/DeepUnits}{52}{57}{Linear.plugAtModuloIso_corolla}

\subsection{Regression proofs: self-loops, mixed depth, and repaired collisions}

For the concrete one-edge self-loop, instantiate shallow associativity with
two unary corollas. Instantiate the right unit with the loop's own corolla as
well. Both retain the loop's shared source/target vertex, demonstrating that
neither theorem relies on a non-self-loop hypothesis.
\lean{Graph/CompositionExamples}{8}{11}{Graph.Examples.loop_associativity}
\lean{Graph/CompositionExamples}{14}{15}{Graph.Examples.loop_unit}

Define a unary variable fragment $V$ as a corolla of the variable label. It has
interface $*\to *$ and no holes: any alleged hole observation would equate
\texttt{none} with \texttt{some}. Let $t$ select the sole edge of the earlier
replacement containing hole one. Its retained inventory is empty, because
there is no other edge in its one-element carrier.
\lean{Graph/CompositionExamples}{22}{23}{Linear.Examples.variableFragment}
\lean{Graph/CompositionExamples}{25}{27}{Linear.Examples.variableFragment_holeFree}
\lean{Graph/CompositionExamples}{30}{30}{Linear.Examples.repeatedSite}
\lean{Graph/CompositionExamples}{32}{34}{Linear.Examples.repeatedSite_retained_absent}

Take the earlier boxed two-hole host and let $s$ select hole zero inside it.
Inserting the one-hole replacement carrying label one gives the already proved
collision with the surviving hole one. This intermediate graph is invalid.
Now fill the inserted occurrence at $s\star t$ with $V$. The composite criterion
proves uniqueness: $V$ is hole-free, hence unique and fresh; and the empty
retained inventory of $t$ is disjoint from the retained host inventory.
\lean{Graph/CompositionExamples}{38}{46}{Linear.Examples.intermediate_collision_repaired}
The final graph still contains the host's hole one, witnessed by the surviving
inner edge. Hole zero is absent: it cannot be retained by $s$ because it was
the selected unique hole; it cannot be retained by $t$ because that inventory
is empty; and it cannot occur in $V$ because $V$ is hole-free. The exact composite
inventory theorem exhausts these three possibilities.
\lean{Graph/CompositionExamples}{49}{60}{Linear.Examples.repaired_inventory}

For a mixed-depth example, form a one-edge shell whose label is the box label
of the boxed two-hole host, and select that shell's sole edge. Insert the boxed
host, then follow its deep first-hole site. Deep associativity compares this
with filling the boxed host first and then inserting it into the shell.
\lean{Graph/CompositionExamples}{63}{64}{Linear.Examples.boxShell}
\lean{Graph/CompositionExamples}{66}{66}{Linear.Examples.shellSite}
\lean{Graph/CompositionExamples}{68}{71}{Linear.Examples.shallow_then_deep_associativity}
A three-stage instance adds the sole site of the repeated-hole replacement and
then the variable fragment; the two site-composition bracketings agree on the
resulting quotient class by three-stage coherence. Finally, restoring the
original fragment at the boxed first-hole site recovers the original boxed
host by deep restoration.
\lean{Graph/CompositionExamples}{74}{79}{Linear.Examples.three_stage_composition}
\lean{Graph/CompositionExamples}{82}{83}{Linear.Examples.deep_restoration}
